\setcounter{aufgabe}{7}
\hypertarget{e2}{}\einheitenkopf{Einheit 2 von 5 · Lineare Funktion $f(x) = m \cdot x + n$}
\zweigzeile{Hier lernst du, Geraden mit $f(x) = m \cdot x + n$ zu zeichnen, $m$ und $n$ abzulesen und zu deuten · neu in diesem Jahr · P10 oft · baut auf: proportionale Funktion (Einheit 1), Brüche, negative Zahlen}

\begin{aufgabe}{\textbf{Ich erkenne $m$ und $n$ in der Gleichung.} Unterstreiche $m$ und kreise $n$ ein. Steht $m$ oder $n$ nicht da, schreibe die fehlende Zahl dazu.}
\begin{teilezwei}
\tz $f(x) = 3x - 5$ & \tz $f(x) = -x + 6$ \\
\tz $f(x) = 7x$ & \tz $f(x) = -4$ \\
\tz $f(x) = 4 + 0{,}5x$ & \\
\end{teilezwei}
\end{aufgabe}

\begin{aufgabe}{\textbf{Ich erkenne, ob eine Gerade steigt oder fällt.} Entscheide ohne Zeichnen: Steigt oder fällt die Gerade?}
\begin{teile}
\teil $f(x) = 2x + 7$ \quad \kreuz{steigt}\kreuz{fällt}
\teil $f(x) = -3x + 1$ \quad \kreuz{steigt}\kreuz{fällt}
\teil $f(x) = 0{,}5x - 4$ \quad \kreuz{steigt}\kreuz{fällt}
\teil $f(x) = -x - 2$ \quad \kreuz{steigt}\kreuz{fällt}
\teil $f(x) = 6 - 2x$ \quad \kreuz{steigt}\kreuz{fällt}
\end{teile}
\end{aufgabe}

\begin{aufgabe}{\textbf{Ich kann ein Steigungsdreieck lesen.} An jeder Geraden ist ein Steigungsdreieck eingezeichnet: erst ein Kästchen nach rechts. Gib an, wie viele Kästchen es dann nach oben oder nach unten geht.}

\begin{ksys}[xmin=-3,xmax=3,ymin=-3,ymax=5,ablesen]
\gerade[2.6]{2}{-2}{p}\gerade[-2.5]{-1}{3}{q}
\draw[very thick] (1,0) -- (2,0) -- (2,2);
\draw[very thick] (-1,4) -- (0,4) -- (0,3);
\end{ksys}\ksysabstand\begin{ksys}[xmin=-3,xmax=3,ymin=-3,ymax=5,ablesen]
\gerade[2.8]{3}{-4}{r}\gerade[-2.6]{-2}{-1}{s}
\draw[very thick] (1,-1) -- (2,-1) -- (2,2);
\draw[very thick] (-2,3) -- (-1,3) -- (-1,1);
\end{ksys}

\begin{teile}
\teil Gerade $p$: \quad \leerfeld{} Kästchen \quad \kreuz{nach oben}\kreuz{nach unten}
\teil Gerade $q$: \quad \leerfeld{} Kästchen \quad \kreuz{nach oben}\kreuz{nach unten}
\teil Gerade $r$: \quad \leerfeld{} Kästchen \quad \kreuz{nach oben}\kreuz{nach unten}
\teil Gerade $s$: \quad \leerfeld{} Kästchen \quad \kreuz{nach oben}\kreuz{nach unten}
\end{teile}
\end{aufgabe}

\begin{aufgabe}{\textbf{Ich kann eine Gerade mit einer Wertetabelle zeichnen.} Fülle die Wertetabelle aus, trage die Punkte ein und zeichne die Gerade. Zeichne alle drei Geraden in dasselbe Koordinatensystem.}
Beispiel: $f(x) = x - 1$ \quad
\wertetabelle[-2,-1,0,1]{x}{f(x)}{-1,0,1,2}
Punkte $(-1\,|\,{-2})$, $(0\,|\,{-1})$, $(1\,|\,0)$, $(2\,|\,1)$ eintragen und mit dem Lineal verbinden.
\begin{teilezwei}
\tz $f(x) = x + 3$ \par \wertetabelle{x}{f(x)}{-1,0,1,2} & \tz $f(x) = 2x - 1$ \par \wertetabelle{x}{f(x)}{-1,0,1,2} \\
\tz $f(x) = -x + 2$ \par \wertetabelle{x}{f(x)}{-1,0,1,2} & \\
\end{teilezwei}

\begin{ksys}[xmin=-2,xmax=3,ymin=-4,ymax=6]
\end{ksys}

\end{aufgabe}

\begin{aufgabe}{\textbf{Ich kann eine Gerade mit $n$ und Steigungsdreieck zeichnen.} Markiere $n$ auf der $y$-Achse. Gehe von dort 1 nach rechts und $m$ nach oben. Zeichne die Gerade durch beide Punkte. Zeichne a) und b) links, c) und d) rechts.}
\beispiel{f(x) = 2x + 3:\quad n &= 3 && \text{Punkt } (0\,|\,3) \\ m &= 2 && \text{1 nach rechts, 2 nach oben: Punkt } (1\,|\,5)}
\begin{teilezwei}
\tz $f(x) = x + 2$ & \tz $f(x) = 3x + 2$ \\
\tz $f(x) = x + 5$ & \tz $f(x) = 2x + 5$ \\
\end{teilezwei}

\begin{ksys}[xmin=-3,xmax=3,ymin=-1,ymax=8]
\end{ksys}\ksysabstand\begin{ksys}[xmin=-3,xmax=3,ymin=-1,ymax=8]
\end{ksys}

\end{aufgabe}

\begin{aufgabe}{\textbf{Ich kann eine Gerade mit $n$ und Steigungsdreieck zeichnen – weiter.} Ist $m$ negativ, geht es nach unten. Zeichne e) und f) links, g) und h) rechts.}
\setcounter{teil}{4}
\begin{teile}
\teil $f(x) = -3x + 5$
\end{teile}
Ist $m$ ein Bruch, gehst du so viele Kästchen nach rechts, wie der Nenner angibt, und so viele nach oben oder unten, wie der Zähler angibt.
\begin{teile}
\teil $f(x) = \frac{1}{2}x + 3$
\teil $f(x) = 2x - 4$
\teil Zeichne den Graphen der Funktion $f$ mit $f(x) = -\frac{2}{3}x - 1$. (P10 2026)
\end{teile}

\begin{ksys}[xmin=-4,xmax=4,ymin=-5,ymax=6]
\end{ksys}\ksysabstand\begin{ksys}[xmin=-4,xmax=4,ymin=-5,ymax=6]
\end{ksys}

\end{aufgabe}

\begin{aufgabe}{\textbf{Ich kann $n$ und $m$ an einer Geraden ablesen.} Lies an jeder Geraden $n$ ab. Zeichne dann ein Steigungsdreieck ein und lies $m$ ab.}

\begin{ksys}[xmin=-4,xmax=5,ymin=-4,ymax=6,ablesen]
\gerade[1.5]{2}{2}{g}\gerade[-1.5]{-1}{4}{h}\gerade[4.5]{0.333333}{-2}{k}
\end{ksys}

\begin{teilezwei}
\tz Gerade $g$: \quad \feld{n} & \tz Gerade $h$: \quad \feld{n} \\
\tz Gerade $k$: \quad \feld{n} & \tz Gerade $g$: \quad \feld{m} \\
\tz Gerade $h$: \quad \feld{m} & \tz Gerade $k$: \quad \feld{m} \\
\end{teilezwei}
\end{aufgabe}

\begin{aufgabe}{\textbf{Ich kann einer Gleichung die passende Gerade zuordnen.} Welche Gerade gehört zur Gleichung? Schreibe ihren Namen auf. Für a) bis d) nimmst du das linke Bild.}

\begin{ksys}[xmin=-3,xmax=3,ymin=-5,ymax=5,ablesen]
\gerade[2]{3}{-2}{p}\gerade[-0.8]{-2}{3}{q}\gerade[2.8]{0.5}{1}{r}\gerade[-2.5]{0}{-3}{s}
\end{ksys}\ksysabstand\begin{ksys}[xmin=-4,xmax=4,ymin=-4,ymax=4,ablesen]
\gerade[-2]{0.5}{2}{l_1}\gerade[3]{-0.5}{-2}{l_2}\gerade[-3]{0.5}{-2}{l_3}\gerade[2.5]{2}{-2}{l_4}
\end{ksys}

\begin{teilezwei}
\tz $y = 3x - 2$: \quad Gerade \leerfeld & \tz $y = -2x + 3$: \quad Gerade \leerfeld \\
\tz $y = 0{,}5x + 1$: \quad Gerade \leerfeld & \tz $y = -3$: \quad Gerade \leerfeld \\
\end{teilezwei}
\begin{teile}
\teil Rechtes Bild: Welche der vier Geraden $l_1$ bis $l_4$ ist der Graph der Funktion mit $y = \frac{1}{2}x - 2$? \quad \leerfeld \quad (P10 2019)
\end{teile}
\end{aufgabe}

\begin{aufgabe}{\textbf{Ich kann an $m$ und $n$ ablesen, wie Geraden zueinander liegen.} Gegeben sind sechs lineare Funktionen:\\
$f_1(x) = 3x - 1$ \quad $f_2(x) = -2x + 4$ \quad $f_3(x) = 3x + 5$ \quad $f_4(x) = 2$ \quad $f_5(x) = -2x - 3$ \quad $f_6(x) = 0{,}5x + 4$}
\begin{teile}
\teil Welche Geraden steigen? \quad \leerfeld
\teil Welche Geraden sind parallel zueinander? \quad \leerfeld
\teil Welche Gerade verläuft parallel zur $x$-Achse? \quad \leerfeld
\teil Welche Geraden schneiden die $y$-Achse im selben Punkt? \quad \leerfeld
\end{teile}
\end{aufgabe}

\begin{aufgabe}{\textbf{Ich finde den Fehler beim Ablesen der Steigung.} Sara liest die Steigung der Geraden $g$ ab. Sie geht von $P$ aus 3 Kästchen nach rechts und 2 Kästchen nach oben zu $Q$ und schreibt:}
\rechnung{m &= \frac{3}{2}}

\begin{ksys}[xmin=-2,xmax=5,ymin=-5,ymax=4,ablesen]
\gerade[4.5]{0.666667}{1}{g}\punkt{0}{1}{P}\punkt{3}{3}{Q}
\gerade[4.5]{-0.5}{-2}{h}
\end{ksys}

\begin{teile}
\teil Finde den Fehler, benenne ihn und gib die richtige Steigung an.
\teil Lies die Steigung der Geraden $h$ ab. \quad \feld{m}
\end{teile}
\end{aufgabe}

\begin{aufgabe}{\textbf{Ich kann begründen, welche Gerade steiler ist.} Entscheide ohne Rechnung und ohne Zeichnen, welche Gerade steiler ist. Begründe.}
\begin{teile}
\teil $g$: $y = 4x - 1$ \quad oder \quad $h$: $y = x + 6$
\teil $g$: $y = -3x + 1$ \quad oder \quad $h$: $y = 2x + 5$
\teil Ben sagt: „$y = 0{,}5x + 8$ ist steiler als $y = 2x$, weil 8 größer ist als 2.“ Stimmt das?
\end{teile}
\end{aufgabe}

\begin{aufgabe}{\textbf{Ich kann $m$ und $n$ in einer Sachsituation deuten.} Zwei volle Regentonnen werden über einen Hahn geleert. Für Tonne A gilt $y = -4x + 120$, für Tonne B gilt $y = -6x + 120$. Dabei ist $x$ die Zeit in Minuten und $y$ die Wassermenge in Litern.}
\begin{teile}
\teil Was bedeuten die Zahlen 120 und $-4$ bei Tonne A im Sachzusammenhang?
\teil Welche Tonne ist zuerst leer? Begründe ohne Rechnung.
\teil Für Tonne C gilt $y = -4x + 90$. Wie liegen die Graphen von A und C zueinander? Was bedeutet das für die beiden Tonnen?
\end{teile}
\end{aufgabe}
